\section{Variance Reduction Techniques for Credit-Linked Notes}

\paragraph{Note on Sources.}
The variance-reduction techniques implemented in this section are drawn from 
\cite{lawKelton2015} \textit{Simulation Modeling and Analysis} (6th ed.), 
specifically Chapter~11.  
The organisation of methods, the logic behind the estimators, 
and much of the conceptual flow follow Law. The notation modified to align with the credit context of this project.

This section introduces three core variance--reduction techniques used to
improve Monte Carlo efficiency when valuing credit--linked notes (CLNs) under
stochastic hazard--rate dynamics: common random numbers, control variates, and
antithetic variates. CLN valuation ultimately reduces to estimating expectations
of functionals of the hazard--rate process. As the complexity of the payoff
increases, particularly for exotic and path--dependent credit products, naive
Monte Carlo becomes inefficient. Variance--reduction techniques are therefore an
essential component of any practical simulation engine.

The techniques discussed here build on the Monte Carlo framework and
hazard--rate simulation machinery developed in sections 2 and 3. Throughout, we
rely on consistent notation, explicit derivations, and examples tailored
specifically to credit products. In particular, the final section illustrates how
each method applies to the pricing of a range--accrual exotic CLN whose
redemption amount depends on the fraction of coupon--observation dates at which
the model--implied CDS spread lies within a predefined corridor.

\subsection{Background on Monte Carlo Estimation}

Monte Carlo methods provide numerical estimates of expectations of the form
\[
\mathbb{E}[H],
\]
where $H$ is a random payoff functional driven by simulated hazard--rate paths
and credit spreads. By drawing independent samples $H_1,\dots,H_N$, the Monte
Carlo estimator is
\[
\hat{H}_N = \frac{1}{N}\sum_{j=1}^N H_j.
\]
The law of large numbers ensures $\hat{H}_N \to \mathbb{E}[H]$ almost surely as
$N\to\infty$.

More importantly, the central limit theorem states that
\[
\sqrt{N}(\hat{H}_N - \mathbb{E}[H])
\]
converges in distribution to a normal with mean zero and variance
$\mathrm{Var}(H)$. The standard error is therefore
\[
SE(\hat{H}_N) = \sqrt{\frac{\mathrm{Var}(H)}{N}}.
\]
This highlights the inefficiency of naive Monte Carlo: halving the standard
error requires quadrupling the number of simulations. Because many credit
payoffs are path--dependent and driven by stochastic term structures,
efficient variance--reduction is critical.

Variance--reduction techniques work by replacing $\hat{H}_N$ with an
alternative estimator that is still unbiased but has a smaller variance. The
remainder of this section develops three such techniques in detail.

\subsection{Common Random Numbers}

Common random numbers (CRNs), also known as correlated sampling, are used when
comparing two system configurations. In our setting, we compare two payoffs
computed under identical simulation noise but different model inputs or payoff
structures. For example, a vanilla CLN and an exotic range--accrual CLN may be
evaluated using the same underlying Gaussian shocks driving the hazard--rate
process.

\subsubsection*{Setup}

Let $X_j^{(1)}$ and $X_j^{(2)}$ be the discounted payoffs for the two
configurations for simulation $j$, constructed using identical random numbers.
Define
\[
Z_j = X_j^{(1)} - X_j^{(2)}.
\]
The quantity of interest is the difference in means
\[
\phi = \mathbb{E}[Z_j] = \mathbb{E}[X^{(1)}] - \mathbb{E}[X^{(2)}].
\]

The CRN estimator is
\[
\tilde{\phi}_N = \frac{1}{N}\sum_{j=1}^N Z_j.
\]
This estimator is unbiased for $\phi$.

\subsubsection*{Variance Derivation}

We compute its variance explicitly:
\[
\mathrm{Var}(\tilde{\phi}_N)
= \frac{1}{N} \mathrm{Var}(Z_j).
\]
Expanding $\mathrm{Var}(Z_j)$:
\[
\mathrm{Var}(Z_j)
= \mathrm{Var}(X_j^{(1)}) + \mathrm{Var}(X_j^{(2)})
- 2\,\mathrm{Cov}(X_j^{(1)}, X_j^{(2)}).
\]

If independent random numbers were used, $\mathrm{Cov}=0$. Under CRNs, the
covariance is typically positive because both payoffs are driven by the same
hazard--rate shocks and payoff structures are similar. A positive covariance
reduces $\mathrm{Var}(Z_j)$ and therefore reduces the variance of
$\tilde{\phi}_N$.

CRNs are heuristic. There is no guarantee that the covariance is positive.
Poorly matched payoffs can increase variance. In practice, CRNs work best when
comparing structurally similar payoffs. The vanilla CLN is a natural candidate
control for the exotic CLN, but its realised effectiveness depends on the chosen
spread corridor and resulting payoff variation.

\subsection{Control Variates}

Control variates (CVs) are among the most powerful variance--reduction
techniques. The method requires a random variable $C$ whose expectation
$\mathbb{E}[C]$ is known analytically and which is strongly correlated with the
target payoff $Y$.

In credit applications, a natural choice is to let
\begin{itemize}
\item $Y$ be the discounted payoff of an exotic CLN,
\item $C$ be the discounted payoff of a vanilla CLN,
\end{itemize}
whose analytical price is available.

\subsubsection*{Derivation}

Consider the adjusted estimator
\[
\tilde{Y}_N = \hat{Y}_N - \beta(\hat{C}_N - \mathbb{E}[C]),
\]
where
\[
\hat{Y}_N = \frac{1}{N}\sum_{j=1}^N Y_j,
\qquad
\hat{C}_N = \frac{1}{N}\sum_{j=1}^N C_j.
\]

This estimator is unbiased regardless of $\beta$. Its variance is
\[
\mathrm{Var}(\tilde{Y}_N)
= \frac{1}{N}\left(
\mathrm{Var}(Y) + \beta^2 \mathrm{Var}(C)
- 2\beta\,\mathrm{Cov}(Y,C)
\right).
\]

Differentiating with respect to $\beta$ gives the optimal coefficient:
\[
\beta^*
= \frac{\mathrm{Cov}(Y,C)}{\mathrm{Var}(C)}.
\]

Substituting back yields
\[
\mathrm{Var}(\tilde{Y}_N)
= \frac{1}{N}
\left[
\mathrm{Var}(Y)\,\bigl(1 - \rho_{Y,C}^2\bigr)
\right],
\]
where $\rho_{Y,C}$ is the correlation between $Y$ and $C$. Variance is reduced
in proportion to the squared correlation.

\subsection{Antithetic Variates}

Antithetic variates use negatively correlated pairs to reduce variance. If $Z$
is a standard normal driving the hazard--rate increment, consider the pair
$(Z, -Z)$. Let the corresponding payoffs be $H(Z)$ and $H(-Z)$.

The antithetic estimator is
\[
\tilde{H}_j = \frac{1}{2}\bigl( H(Z_j) + H(-Z_j) \bigr),
\qquad
\hat{H}_N = \frac{1}{N}\sum_{j=1}^N \tilde{H}_j.
\]

Explicitly computing the variance:
\[
\mathrm{Var}(\tilde{H}_j)
= \frac{1}{4}\left(
\mathrm{Var}(H(Z)) + \mathrm{Var}(H(-Z))
+ 2\,\mathrm{Cov}(H(Z),H(-Z))
\right).
\]
Because $H(Z)$ and $H(-Z)$ often move in opposite directions for monotone
payoffs, the covariance is frequently negative, reducing variance. As with
CRNs, this is not guaranteed; highly nonlinear payoffs may yield little or no
benefit.

\subsection{Application to Exotic CLNs}

We now apply these techniques to the exotic range--accrual CLN whose payoff
depends on how often the model--implied CDS spread lies within a corridor
$[L,U]$ over the coupon observation dates.

\paragraph{Control Variates.}
Let $Y$ be the discounted payoff of the exotic CLN, and let $C$ be the
discounted payoff of a vanilla CLN with the same notional, maturity, and
recovery rate. Because the exotic structure modifies only the redemption while
leaving the coupon and recovery legs close to the vanilla structure, this is a
plausible control variate. Whether it produces strong variance reduction is an
empirical question governed by the realised correlation between $Y$ and $C$.

\paragraph{Enhancing Control Variates Using Common Random Numbers.}
A further improvement is obtained by coupling $Y$ and $C$ via identical random
numbers. This ensures that the two payoffs react to the same hazard--rate
shocks and CDS--spread paths. CRNs increase the correlation between $Y$ and $C$
and therefore increase the effectiveness of the control variate.

\paragraph{Antithetic Variates.}
Antithetic variates can be applied at the driver level: for each simulation
$j$, generate hazard--rate shocks using $Z_j$ and $-Z_j$. Compute exotic CLN
payoffs for both and average them. In the next section we will see if the negatively
correlated pair of random variables leads to a meaningful variance reduction. 
